% !TeX root = ../complete.tex
\section{برآورد \lr{Dynamic Power}}
\label{sec:power-dynamic}

\subsection{از شمارش گذار تا توان سوئیچینگ}
\label{sec:power-switching}

در یک بازه‌ی مشاهده به طول \(T\)، تعداد شارژهای کامل یک گره از صفر
تا \(V_{DD}\) را \(N_{\uparrow}\) می‌نامیم. با فرض‌های بخش
\ref{sec:power-capacitor-energy}، هر شارژ انرژی \(C_LV_{DD}^2\) را
از منبع می‌گیرد. \(f_{\mathrm{sw}}=N_{\uparrow}/T\) نرخ شارژها
برحسب تعداد در ثانیه است. بنابراین توان متوسطی که منبع برای شارژ بار
تحویل می‌دهد از رابطه زیر به دست می‌آید
\cite[رابطه‌های \secnum{۵.۹} و \secnum{۵.۱۰}، صفحه ۱۸۴]{weste}
\cite[اسلایدهای ۷ و ۸]{powerslides}.
\begin{equation}
  P_{\mathrm{sw}}=\frac{N_{\uparrow}C_LV_{DD}^2}{T}
                =C_LV_{DD}^2 f_{\mathrm{sw}}.
  \label{eq:power-switching-rate}
\end{equation}
برابری انرژی گرفته‌شده از منبع و انرژی تلف‌شده برای یک چرخه‌ی کامل
برقرار است. در بازه‌ای که ولتاژ ابتدا و انتهای آن متفاوت باشد، برای
حساب انرژی تلف‌شده باید تغییر انرژی ذخیره‌شده در خازن را هم لحاظ کرد.

\lr{Clock} سیگنالی دوره‌ای است که زمان عمل مدار را تعیین می‌کند.
اگر فرکانس ساعت \(f\) باشد، در بازه‌ای به طول \(T\) که دقیقاً شامل
تعدادی دوره کامل است، \(N_{\mathrm{clk}}=fT\) دوره داریم. با تقسیم
تعداد شارژها بر تعداد دوره‌ها، \lr{Activity Factor} را تعریف می‌کنیم.
\begin{equation}
  \alpha=\frac{N_{\uparrow}}{N_{\mathrm{clk}}}
        =\frac{f_{\mathrm{sw}}}{f}.
  \label{eq:power-activity-count}
\end{equation}
صورت و مخرجِ کسر اول شمارش‌اند. پس \(\alpha\) بی‌بعد است و میانگین
تعداد گذار صعودی در هر دوره ساعت را نشان می‌دهد. با جای‌گذاری
\(f_{\mathrm{sw}}=\alpha f\) در رابطه \eqref{eq:power-switching-rate}،
برآورد توان سوئیچینگ برحسب فرکانس ساعت چنین می‌شود.
\begin{equation}
  P_{\mathrm{sw}}=\alpha C_LV_{DD}^2 f.
  \label{eq:power-switching}
\end{equation}
در این قرارداد، تنها گذار \(0\to 1\) شمرده می‌شود. خود ساعت در هر
دوره یک لبه صعودی و یک لبه نزولی دارد، پس \(\alpha=1\) است. برای
داده‌ای که در هر دوره یک بار مقدارش عوض می‌شود، گذارها یکی‌درمیان
صعودی و نزولی‌اند. در یک بازه شامل تعداد زوج از این گذارها،
\(\alpha=1/2\) است. گره ثابت \(\alpha=0\) دارد. اگر خود گره بار
اسلاید ۶ در هر دوره ساعت \qty{1}{\giga\hertz} یک بار شارژ شود،
از انرژی بخش \ref{sec:power-cycle-example} به توان زیر می‌رسیم.
\[
  P_{\mathrm{sw}}=(\qty{150e-15}{\joule})(\qty{1e9}{\hertz})
                =\qty{150}{\micro\watt}.
\]
هرتز معادل معکوس ثانیه است. اگر همین گره به‌طور متوسط در هر دو دوره
یک بار شارژ شود، با همان انرژی به‌ازای هر شارژ، توان سوئیچینگ نصف
می‌شود.

\subsection{جمع سهم گره‌ها و تخمین توان یک تراشه}
\label{sec:power-effective-capacitance}

در یک شبکه چندگرهی، هر گره شماره \(i\) بار معادل \(C_i\) و
\lr{Activity Factor} مختص خود \(\alpha_i\) را دارد. بار شامل
ظرفیت ورودی گیت‌های متصل به گره، ظرفیت‌های \lr{Diffusion} روی گره
و ظرفیت سیم است. در مدل ساده، ولتاژ تغذیه و فرکانس را برای همه‌ی گره‌ها
مشترک می‌گیریم و \(\alpha_i C_i\) را ظرفیت مؤثرِ فعالیت آن گره
می‌نامیم. برای محاسبه توان، سهم گره‌ها را جمع می‌کنیم
\cite[بخش \secnum{۵.۲}، صفحه ۱۸۵]{weste}.
\begin{equation}
  P_{\mathrm{sw}}=V_{DD}^2 f\sum_i\alpha_i C_i.
  \label{eq:power-node-sum}
\end{equation}
برای گره‌هایی که تغذیه یا فرکانس متفاوت دارند، توان هر گروه را با
مقادیر همان گروه حساب می‌کنیم و سپس توان‌ها را با هم جمع می‌کنیم.

در مثال ۵.۱ کتاب و اسلایدهای ۱۱ و ۱۲، مدار به دو گروه منطق و حافظه
تقسیم شده است. در هر گروه، به‌جای محاسبه تک‌تک گره‌ها، از عرض متوسط
ترانزیستورها و یک \lr{Activity Factor} نماینده استفاده می‌کنیم. تعداد
ترانزیستورهای منطق \(50\times10^6\) و تعداد ترانزیستورهای حافظه
\(950\times10^6\) است. اگر واحد هندسی \(\lambda\) برابر
\qty{25}{\nano\metre} باشد، عرض‌های متوسط \(12\lambda\) و
\(4\lambda\) را از راه تبدیل واحد به‌ترتیب
\qty{0.30}{\micro\metre} و \qty{0.10}{\micro\metre} به دست
می‌آوریم. فرض مثال برای مجموع ظرفیت \lr{Gate} و \lr{Diffusion}،
\qty{1.8}{\femto\farad\per\micro\metre} به‌ازای هر واحد عرض است.
ظرفیت سیم و توان جریان مستقیم در این تخمین کنار گذاشته می‌شوند
\cite[مثال ۵.۱، صفحات ۱۸۵ و ۱۸۶]{weste}
\cite[اسلایدهای ۱۱ و ۱۲]{powerslides}.

با ضرب تعداد ترانزیستور در عرض متوسط و ظرفیت بر واحد عرض، ظرفیت
تجمیعی هر گروه به دست می‌آید. واحد طول در ضرب حذف می‌شود و حاصل
برحسب فاراد است.
\[
  \begin{aligned}
    C_{\mathrm{logic}}&=(50\times10^6)(\qty{0.30}{\micro\metre})
      (\qty[per-mode=symbol]{1.8}{\femto\farad\per\micro\metre})
      =\qty{27}{\nano\farad},\\
    C_{\mathrm{mem}}&=(950\times10^6)(\qty{0.10}{\micro\metre})
      (\qty[per-mode=symbol]{1.8}{\femto\farad\per\micro\metre})
      =\qty{171}{\nano\farad}.
  \end{aligned}
\]
با فعالیت \(0.1\) برای منطق و \(0.02\) برای حافظه، ظرفیت‌های مؤثر
\qty{2.7}{\nano\farad} و \qty{3.42}{\nano\farad} هستند. پس با
تغذیه \qty{1.0}{\volt} و ساعت \qty{1.0}{\giga\hertz}، از رابطه
\eqref{eq:power-node-sum} به برآورد زیر می‌رسیم.
\[
  P_{\mathrm{sw}}=(\qty{2.7}{\nano\farad}+\qty{3.42}{\nano\farad})
                 \times(\qty{1.0}{\volt})^2\times\qty{1.0}{\giga\hertz}
                 \approx\qty{6.1}{\watt}.
\]
در این تخمین، ظرفیت کل حافظه
بیشتر است ولی سهم آن از توان به همان نسبت بزرگ نمی‌شود، چون فعالیت
کمتری دارد. برای استفاده از این روش در یک تراشه دیگر، فعالیت باید
نماینده‌ی بار کاری همان تراشه باشد.

\subsection{جریان مستقیم در گذار و مؤلفه‌های توان}
\label{sec:power-short-circuit}

در یک گذار واقعی، ورودی \lr{Inverter} مدتی بین صفر و \(V_{DD}\)
می‌ماند. در مدل آستانه‌ای، \lr{nMOS} با \(v_{\mathrm{in}}>V_{tn}\)
روشن می‌شود و \lr{pMOS} تا \(v_{\mathrm{in}}<V_{DD}-|V_{tp}|\)
روشن می‌ماند. \(V_{tn}\) آستانه \lr{nMOS} و \(|V_{tp}|\) بزرگی
آستانه \lr{pMOS} است. بنابراین اگر \(V_{tn}+|V_{tp}|<V_{DD}\)،
در بازه‌ی بین دو آستانه هر دو ترانزیستور می‌توانند هم‌زمان هدایت کنند.
بخشی از جریان مستقیماً از تغذیه به زمین می‌رود و در شارژ خازن سهمی
ندارد. این سهم \lr{Short-Circuit Current} نام دارد. با ورودی کند،
مدت هدایت مشترک می‌تواند بیشتر شود. مقدار انرژی به شکل موج ورودی
و خروجی و بار گره وابسته است \cite[بخش \secnum{۵.۲.۵}، صفحه ۱۹۳]{weste}
\cite[اسلاید ۹]{powerslides}.

برای گیتی با ابعاد و بار مشخص، مجموع انرژیِ مربوط به این پالس‌های
جریان را بر زمان مشاهده تقسیم می‌کنیم تا توان \(P_{\mathrm{sc}}\)
به دست آید.
اسلاید ۹ در شرایط قابل‌مقایسه بودن زمان‌های گذار ورودی و خروجی، این
سهم را کوچک برآورد می‌کند. برای برآورد دستی در آن شرایط می‌توان از
\(P_{\mathrm{dyn}}\approx P_{\mathrm{sw}}\) استفاده کرد. در جمع
توان، \(P_{\mathrm{stat}}\) مؤلفه‌ای است که با ثابت ماندن وضعیت منطقی
نیز می‌تواند وجود داشته باشد. رابطه جمع مؤلفه‌ها چنین است
\cite[بخش \secnum{۵.۱.۳}، صفحات ۱۸۴ و ۱۸۵]{weste}
\cite[اسلاید ۱۰]{powerslides}.
\begin{equation}
  P_{\mathrm{dyn}}=P_{\mathrm{sw}}+P_{\mathrm{sc}},\qquad
  P_{\mathrm{total}}=P_{\mathrm{dyn}}+P_{\mathrm{stat}}.
  \label{eq:power-components}
\end{equation}
اگر گره‌ها گذار نداشته باشند، سهم سوئیچینگ و پالس جریان مستقیم ناشی
از گذار صفر می‌شود. نشتی‌های معرفی‌شده در
\NotesTopic{sec:leakage}{بحث جریان‌های نشتی} از مسیرهای باقی‌مانده
برای مصرف توان‌اند.
